\chapter{要素的表述}
\section{图}
图与其前面的条文之间空一行，图题位于图下方。图中的注位于图脚注之上，图脚注位于图题之上。\index{图}

\begin{figure}[htb]
    \stdunits{单位为毫米}
    \centering
    \rule{4cm}{3cm}\textsuperscript{a}
    \stdnote{这是图中的注。}
    \stdfootnotetext{a}{这是图脚注。}
    \caption{图题示例}
    \label{fig:example}
\end{figure}

\section{表}
表题位于表上方，表中的注位于表内下方、表脚注之上。\index{表}

\begin{table}[htb]
    \stdunits{单位为毫米}
    \caption{表题示例}
    \label{tab:example}
    \centering\zihao{-5}\songti
    \begin{tabular}{|c|c|}
        \hline
        项目 & 长度 \\
        \hline
        示例甲 & 10 \\
        示例乙 & 20\textsuperscript{a} \\
        \hline
    \end{tabular}
    \stdnote{这是表中的注。}
    \stdfootnotetext{a}{这是表脚注。}
\end{table}

\subsection{表的转页接排}
续表重复表编号和表头，并标明“续”。\index{续表}
{\zihao{-5}\songti
\setlength{\LTleft}{\fill}
\setlength{\LTright}{\fill}
\begin{longtable}{|c|c|}
    \caption{跨页表题示例}
    \label{tab:continued}\\
    \hline
    序号 & 内容 \\
    \hline
    \endfirsthead
    \caption*{\thetable\quad（续）}\\
    \hline
    序号 & 内容 \\
    \hline
    \endhead
    1 & 内容 01 \\
    2 & 内容 02 \\
    3 & 内容 03 \\
    4 & 内容 04 \\
    5 & 内容 05 \\
    6 & 内容 06 \\
    7 & 内容 07 \\
    8 & 内容 08 \\
    9 & 内容 09 \\
    10 & 内容 10 \\
    11 & 内容 11 \\
    12 & 内容 12 \\
    13 & 内容 13 \\
    14 & 内容 14 \\
    15 & 内容 15 \\
    16 & 内容 16 \\
    17 & 内容 17 \\
    18 & 内容 18 \\
    19 & 内容 19 \\
    20 & 内容 20 \\
    21 & 内容 21 \\
    22 & 内容 22 \\
    23 & 内容 23 \\
    24 & 内容 24 \\
    25 & 内容 25 \\
    26 & 内容 26 \\
    27 & 内容 27 \\
    28 & 内容 28 \\
    29 & 内容 29 \\
    30 & 内容 30 \\
    31 & 内容 31 \\
    32 & 内容 32 \\
    33 & 内容 33 \\
    34 & 内容 34 \\
    35 & 内容 35 \\
    36 & 内容 36 \\
    37 & 内容 37 \\
    38 & 内容 38 \\
    39 & 内容 39 \\
    40 & 内容 40 \\
    41 & 内容 41 \\
    42 & 内容 42 \\
    43 & 内容 43 \\
    44 & 内容 44 \\
    45 & 内容 45 \\
    46 & 内容 46 \\
    47 & 内容 47 \\
    48 & 内容 48 \\
    49 & 内容 49 \\
    50 & 内容 50 \\
    51 & 内容 51 \\
    52 & 内容 52 \\
    53 & 内容 53 \\
    54 & 内容 54 \\
    55 & 内容 55 \\
    56 & 内容 56 \\
    57 & 内容 57 \\
    58 & 内容 58 \\
    59 & 内容 59 \\
    60 & 内容 60 \\
    \multicolumn{2}{l}{注：跨页表结束时给出表注。}\\
\end{longtable}
}

\section{数学公式}
数学公式另行居中编排，编号右端对齐。公式编号从引言开始连续到附录之前。\index{数学公式}

\begin{equation}
    E = mc^2
    \label{eq:energy}
\end{equation}

式中：

$E$——能量；

$m$——质量；

$c$——真空中的光速。

\section{注、示例和脚注}
注只给出有助于理解或使用文件内容的附加信息，不包含要求型条款。\index{注}
\stdnote{这是条文中的单个注。}
\stdnote[1]{这是条文中的第一个注。}
\stdnote[2]{这是条文中的第二个注。}

示例用于帮助理解条款。示例不宜单独设章或条。\index{示例}
\stdexample{这是单个示例。}
\stdexample[1]{这是第一个示例。}
\stdexample[2]{这是第二个示例。}
\stdexampleframe{这是置于线框内的示例，可用于展示编排格式或避免与条款混淆。}

正文脚注使用带半圆括号的阿拉伯数字编号\footnote{这是条文脚注，编号、缩进和分隔线均按规范设置。}。

\section{引用和提示}
\ref{fig:example}、\ref{tab:example}、\ref{tab:continued} 和 \eqref{eq:energy} 分别展示图、表和公式的引用。\index{交叉引用}

附录~\ref{appendix:normative} 和附录~\ref{appendix:informative} 展示规范性及资料性附录的提示。资料性引用可使用文献条目\cite{lamport1994latex}。
